\chapter{Peubah Acak dan Distribusi Binomial}\label{ch:g12:randvar}

\omterm{def:g12:randvar:rv}{Peubah acak} melekatkan sebuah bilangan pada tiap hasil percobaan: sebuah
keuntungan, sebuah cacah, sebuah lama waktu. Harapannya adalah rata-rata jangka
panjang nilai yang dihasilkannya, dan \omterm{def:g12:randvar:exp}{ragamnya} mengukur penyebarannya. Bintang bab
ini adalah \omterm{def:g12:randvar:binomial}{distribusi binomial}, yang mencacah keberhasilan pada ulangan percobaan
yang \omterm{def:g12:condprob:indep}{saling bebas}. \omterm{def:g12:randvar:rv}{Peubah acak} dan \omterm{def:g12:randvar:binomial}{distribusi binomial} pertama kali
dijumpai di \cref{ch:g11:prob,ch:g11:binom}; bab ini meninjau kembali dan
memperdalamnya, dengan perkakas kombinatorika \cref{ch:g12:comb} yang kini
tersedia.

\section{Peubah acak diskret}

\begin{definition}[Peubah acak, distribusi]\label{def:g12:randvar:rv}
Suatu \emph{peubah acak}\index{peubah acak} pada \omterm{def:g11:prob:model}{ruang sampel} berhingga
$\Omega$ adalah \omterm{def:g11:func:function}{fungsi} $X \colon \Omega \to \R$. Yang disebut
\emph{distribusinya}\index{distribusi} (atau hukumnya) adalah data nilai
yang mungkin $x_1, \dots, x_k$ beserta peluangnya
\[
p_i = \P(X = x_i), \qquad \sum_{i=1}^{k} p_i = 1 .
\]
\end{definition}

\begin{definition}[Nilai harapan, ragam, simpangan baku]\label{def:g12:randvar:exp}
\emph{Nilai harapan}\index{nilai harapan} dari $X$ adalah
\[
\E(X) = \sum_{i=1}^{k} p_i\, x_i ,
\]
sedangkan \emph{ragamnya}\index{ragam} dan \emph{\omterm{def:g11:stat:variance}{simpangan bakunya}} adalah
\[
\V(X) = \E\bigl((X - \E(X))^2\bigr) = \sum_{i=1}^k p_i\bigl(x_i - \E(X)\bigr)^2,
\qquad
\sigma(X) = \sqrt{\V(X)} .
\]
\end{definition}

\begin{proposition}[Rumus König--Huygens]\label{prop:g12:randvar:konig}
\index{rumus König--Huygens}
$\V(X) = \E(X^2) - \E(X)^2$.
\end{proposition}

\begin{proof}
Tulislah $m = \E(X)$ lalu jabarkan:
\[
\V(X) = \sum_i p_i (x_i^2 - 2m x_i + m^2)
= \E(X^2) - 2m\sum_i p_i x_i + m^2\sum_i p_i
= \E(X^2) - 2m^2 + m^2 . \qedhere
\]
\end{proof}

\begin{proposition}[Transformasi afin]\label{prop:g12:randvar:affine}
Untuk $a, b \in \R$:
\[
\E(aX + b) = a\,\E(X) + b, \qquad
\V(aX + b) = a^2\,\V(X) .
\]
\end{proposition}

\begin{proof}
Yang pertama hanyalah penyusunan ulang jumlah pendefinisinya. Untuk yang kedua, $aX + b$
menyimpang dari rata-ratanya sebesar $a(X - \E(X))$, dan pengkuadratannya mengalikan
dengan $a^2$.
\end{proof}

\begin{example}[Permainan yang adil]\label{ex:g12:randvar:game}
Sebuah permainan berbiaya $2$ euro; sebuah dadu dilempar, dan pemainnya menerima nilai
yang muncul jika nilainya sekurang-kurangnya $5$, dan tidak menerima apa-apa jika tidak. Keuntungannya $G$ bernilai
$-2$ (peluangnya $\frac46$), $3$ ($\frac16$), $4$ ($\frac16$):
\[
\E(G) = \frac{-8 + 3 + 4}{6} = -\frac{1}{6} < 0 .
\]
\omterm{def:g11:stat:mean}{Rata-rata} pemainnya rugi $17$ sen tiap permainan: jadi permainannya merugikan
(seperti kebanyakan permainan yang nyata).
\end{example}

\section{Percobaan Bernoulli dan distribusi binomial}

\begin{definition}[Distribusi Bernoulli]\label{def:g12:randvar:bernoulli}
Suatu \emph{percobaan Bernoulli}\index{percobaan Bernoulli} adalah percobaan dengan dua
hasil, \emph{berhasil} (peluangnya $p$) dan \emph{gagal}
($q = 1 - p$). Penunjuk $X$ bagi keberhasilan ($X = 1$ jika berhasil, $0$ jika
gagal) mengikuti \emph{\omterm{def:g12:randvar:rv}{distribusi} Bernoulli} $\mathcal B(p)$:
\[
\E(X) = p, \qquad \V(X) = p(1-p) .
\]
(Memang $\E(X) = p$, $\E(X^2) = p$, dan König--Huygens memberi
$\V(X) = p - p^2$.)
\end{definition}

\begin{definition}[Distribusi binomial]\label{def:g12:randvar:binomial}
Ulangilah \omterm{def:g12:randvar:bernoulli}{percobaan Bernoulli} $n$ kali secara \omterm{def:g12:condprob:indep}{saling bebas}, lalu misalkan $X$ cacah
total keberhasilannya. \omterm{def:g12:randvar:rv}{Distribusi} $X$ itulah yang disebut \emph{distribusi
binomial}\index{distribusi binomial} $\mathcal B(n, p)$.
\end{definition}

\begin{omfigure}
\begin{tikzpicture}[font=\small]
  \coordinate (root) at (0,0);
  \node (S) at (1.6,1.6) {B};
  \node (F) at (1.6,-1.6) {G};
  \node (SS) at (3.2,2.4) {B};
  \node (SF) at (3.2,0.8) {G};
  \node (FS) at (3.2,-0.8) {B};
  \node (FF) at (3.2,-2.4) {G};
  \node (SSS) at (4.8,2.8) {B};
  \node (SSF) at (4.8,2.0) {G};
  \node (SFS) at (4.8,1.2) {B};
  \node (SFF) at (4.8,0.4) {G};
  \node (FSS) at (4.8,-0.4) {B};
  \node (FSF) at (4.8,-1.2) {G};
  \node (FFS) at (4.8,-2.0) {B};
  \node (FFF) at (4.8,-2.8) {G};
  \draw[omThm, thick] (root) -- (S);
  \draw[omThm, thick] (S) -- (SS);   \draw[omThm, thick] (SS) -- (SSF);
  \draw[omThm, thick] (S) -- (SF);   \draw[omThm, thick] (SF) -- (SFS);
  \draw[omThm, thick] (root) -- (F);
  \draw[omThm, thick] (F) -- (FS);   \draw[omThm, thick] (FS) -- (FSS);
  \draw[gray] (SS) -- (SSS);
  \draw[gray] (SF) -- (SFF);
  \draw[gray] (FS) -- (FSF);
  \draw[gray] (F) -- (FF);
  \draw[gray] (FF) -- (FFS);
  \draw[gray] (FF) -- (FFF);
  \node[gray, anchor=west] at (5.3,2.8) {$X=3$};
  \node[omThm, anchor=west] at (5.3,2.0) {$X=2$};
  \node[omThm, anchor=west] at (5.3,1.2) {$X=2$};
  \node[gray, anchor=west] at (5.3,0.4) {$X=1$};
  \node[omThm, anchor=west] at (5.3,-0.4) {$X=2$};
  \node[gray, anchor=west] at (5.3,-1.2) {$X=1$};
  \node[gray, anchor=west] at (5.3,-2.0) {$X=1$};
  \node[gray, anchor=west] at (5.3,-2.8) {$X=0$};
\end{tikzpicture}

{\small Tiga \omterm{def:g12:randvar:bernoulli}{percobaan Bernoulli}: tepat $\binom{3}{2} = 3$ dari
$2^3$ lintasannya memberi dua keberhasilan (merah), masing-masing berpeluang $p^2(1-p)$,
sehingga $\P(X = 2) = 3p^2(1-p)$.}
\end{omfigure}

\begin{theorem}\label{thm:g12:randvar:binomial}
Jika $X \sim \mathcal B(n, p)$, maka untuk $0 \leq k \leq n$:
\[
\P(X = k) = \binom{n}{k} p^k (1-p)^{n-k},
\]
dan
\[
\E(X) = np, \qquad \V(X) = np(1-p) .
\]
\end{theorem}

\begin{proof}
Suatu \omterm{def:g12:seq:sequence}{barisan} hasil tertentu dengan $k$ keberhasilan dan $n-k$ kegagalan
berpeluang $p^k(1-p)^{n-k}$ menurut \omterm{def:g12:condprob:indep}{kebebasannya}; sedangkan banyaknya \omterm{def:g12:seq:sequence}{barisan} semacam itu
adalah banyaknya cara menempatkan $k$ keberhasilan di antara $n$ percobaannya,
yaitu $\binom nk$ (\cref{ch:g12:comb}). Menjumlahkan atas \omterm{def:g12:seq:sequence}{barisannya} memberi
rumusnya --- dan teorema binomial menegaskan
$\sum_k \P(X = k) = (p + q)^n = 1$.

Untuk harapannya, dengan memakai $k\binom nk = n\binom{n-1}{k-1}$
(\cref{exo:g12:comb:7}):
\[
\E(X) = \sum_{k=1}^{n} k\binom nk p^k q^{n-k}
= np\sum_{k=1}^{n} \binom{n-1}{k-1} p^{k-1} q^{(n-1)-(k-1)}
= np\,(p + q)^{n-1} = np .
\]
Rumus \omterm{def:g12:randvar:exp}{ragamnya} dibuktikan serupa dengan identitas
$k(k-1)\binom nk = n(n-1)\binom{n-2}{k-2}$, yang memberi
$\E(X(X-1)) = n(n-1)p^2$, sehingga
$\V(X) = n(n-1)p^2 + np - (np)^2 = np(1-p)$. (Bukti yang lebih berstruktur --- yaitu
\omterm{def:g12:randvar:exp}{ragam} jumlah peubah yang \omterm{def:g12:condprob:indep}{saling bebas} --- datang bersama
\cref{thm:g12:sums:variance}.)
\end{proof}

\begin{omfigure}
\begin{tikzpicture}
\begin{axis}[
  width=10cm, height=4.8cm,
  ybar, bar width=4pt,
  xlabel={$k$}, ylabel={$\P(X = k)$},
  xmin=-1, xmax=21, ymin=0,
  ytick={0, 0.1, 0.2},
  title={$\mathcal B(20,\ 0.3)$}]
  \addplot[fill=omDef!40, draw=omDef] coordinates {
    (0,0.000798) (1,0.006839) (2,0.027846) (3,0.071604)
    (4,0.130421) (5,0.178863) (6,0.191639) (7,0.164262)
    (8,0.114397) (9,0.065370) (10,0.030817) (11,0.012006)
    (12,0.003859) (13,0.001018) (14,0.000218) (15,0.000037)
    (16,0.000005) (17,0.000001) (18,0) (19,0) (20,0)};
\end{axis}
\end{tikzpicture}

{\small \omterm{def:g12:randvar:rv}{Distribusi} $\mathcal B(20, 0.3)$: rata-ratanya $np = 6$, \omterm{def:g11:stat:variance}{simpangan
bakunya} $\sqrt{npq} \approx 2.05$.}
\end{omfigure}

\begin{method}[Mengenali keadaan binomial]\label{met:g12:randvar:recognize}
Periksalah ketiga bahannya sebelum menulis $X \sim \mathcal B(n,p)$: yaitu
\emph{cacah yang tetap} $n$ bagi percobaannya; \emph{dua hasil} tiap percobaannya dengan
peluang berhasil $p$ yang \emph{sama}; serta \emph{\omterm{def:g12:condprob:indep}{kebebasan}} percobaannya
(pengambilan contoh \emph{dengan} pengembalian, atau dari populasi yang besar). Lalu pakailah
\[
\P(X \geq 1) = 1 - (1-p)^n
\]
untuk ``sekurang-kurangnya satu keberhasilan'', dan kalkulator atau tabel kumulatif untuk
$\P(X \leq k)$ yang umum.
\end{method}

\begin{example}\label{ex:g12:randvar:atleastone}
Berapa kali orang harus melempar dadu agar berpeluang sekurang-kurangnya $99\%$
memperoleh angka enam? Dengan $X \sim \mathcal B\left(n, \frac16\right)$:
$\P(X \geq 1) = 1 - \left(\frac56\right)^n \geq 0.99$ berarti
$\left(\frac56\right)^n \leq 0.01$, \emph{yaitu}
$n \geq \frac{\ln 0.01}{\ln(5/6)} \approx 25.3$: jadi mulai dari $n = 26$ lemparan.
\end{example}

\section{Latihan}

\begin{exercise}[$\star$]\label{exo:g12:randvar:1}
Suatu \omterm{def:g12:randvar:rv}{peubah acak} $X$ bernilai $-1, 0, 2, 5$ dengan peluang
$0.3, 0.2, 0.4, 0.1$. Hitunglah $\E(X)$, $\V(X)$ dan $\sigma(X)$.
\end{exercise}

\begin{exercise}[$\star$]\label{exo:g12:randvar:2}
Sebuah uji pilihan ganda berisi $10$ pertanyaan, masing-masing dengan $4$ pilihan, dan satu
di antaranya benar. Seorang murid menjawab secara acak seragam dan \omterm{def:g12:condprob:indep}{saling bebas}.
Misalkan $X$ cacah jawaban yang benar.
\begin{enumerate}
  \item Berikan \omterm{def:g12:randvar:rv}{distribusi} $X$, lalu $\E(X)$ dan $\sigma(X)$.
  \item Hitunglah $\P(X = 0)$, $\P(X = 5)$ dan $\P(X \geq 1)$.
\end{enumerate}
\end{exercise}

\begin{exercise}[$\star$]\label{exo:g12:randvar:3}
Sebuah perusahaan asuransi menanggung $n = 400$ nasabah; masing-masing mengajukan klaim selama
setahun dengan peluang $p = 0.05$, secara \omterm{def:g12:condprob:indep}{saling bebas}. Misalkan $X$ cacah
klaimnya. Kenalilah \omterm{def:g12:randvar:rv}{distribusi} $X$ lalu hitung harapannya dan
\omterm{def:g11:stat:variance}{simpangan bakunya}.
\end{exercise}

\begin{exercise}[$\star\star$]\label{exo:g12:randvar:4}
Pada permainan \cref{ex:g12:randvar:game}, penyelenggaranya menginginkan permainan yang adil
($\E(G) = 0$) dengan mengubah harga masuknya $c$. Carilah $c$. Hitunglah
\omterm{def:g12:randvar:exp}{ragam} keuntungannya bagi versi yang adil ini; apakah ``adil'' sama dengan
``tanpa risiko''?
\end{exercise}

\begin{exercise}[$\star\star$]\label{exo:g12:randvar:5}
Seorang pemain bola basket mencetak lemparan bebas dengan peluang $0.7$. Ia menembak
$8$ kali (tembakan yang \omterm{def:g12:condprob:indep}{saling bebas}). Hitunglah peluang bahwa ia mencetak:
tepat $6$; sekurang-kurangnya $6$; sekurang-kurangnya satu. Berapakah cacah cetakan yang
paling mungkin?
\end{exercise}

\begin{exercise}[$\star\star$]\label{exo:g12:randvar:6}
Sebuah maskapai penerbangan tahu bahwa tiap penumpang yang memesan datang dengan peluang
$0.9$, secara \omterm{def:g12:condprob:indep}{saling bebas}. Sebuah penerbangan berisi $100$ kursi dan maskapainya menjual $104$
tiket. Nyatakan, dengan memakai \omterm{def:g12:randvar:binomial}{distribusi binomial}, peluang bahwa penumpang
yang datang lebih banyak daripada kursinya, lalu batasi secara numeris dengan
kalkulator (berikan ungkapan eksaknya).
\end{exercise}

\begin{exercise}[$\star\star$]\label{exo:g12:randvar:7}
Misalkan $X \sim \mathcal B(n, p)$. Tunjukkan bahwa
\[
\frac{\P(X = k+1)}{\P(X = k)} = \frac{n-k}{k+1}\cdot\frac{p}{1-p},
\]
lalu simpulkan bahwa \omterm{def:g12:randvar:rv}{distribusinya} naik sampai
$k^* = \floor{(n+1)p}$ dan turun sesudahnya (itulah \emph{modus} \omterm{def:g12:randvar:binomial}{distribusi
binomialnya}).
binomial).
\end{exercise}

\begin{exercise}[$\star\star\star$]\label{exo:g12:randvar:8}
\emph{(Saint Petersburg yang dijinakkan.)} Sekeping uang logam yang setimbang dilempar sampai gambar
muncul, tetapi paling banyak $10$ kali. Misalkan $N$ cacah lemparan yang terpakai, dan
pemainnya menerima $2^N$ euro jika gambar muncul, dan $0$ jika tidak.
\begin{enumerate}
  \item Berikan \omterm{def:g12:randvar:rv}{distribusi} $N$ yang dibatasi pada hasil yang menang:
        $\P(\text{gambar pertama pada lemparan } k) = 2^{-k}$ untuk
        $1 \leq k \leq 10$, lalu periksalah peluang total kemenangannya.
  \item Hitunglah harapan bayarannya. Akan menjadi berapa tanpa batas
        $10$ lemparan itu?
\end{enumerate}
\end{exercise}

\section{Soal: Penerbangan yang kelebihan pesanan}

\begin{problem}[{Soal akhir pekan --- maskapai menjual kursi lebih banyak daripada
yang dimilikinya, pabrik menerima lot yang nyaris tak diperiksanya, dan
distribusi binomial mewasiti keduanya}]
\label{pb:g12:randvar:1}
Sebuah maskapai berkursi $100$ dengan riang menjual $105$ tiket: kira-kira
$10\,\%$ penumpangnya tak pernah datang, dan kursi kosong tak menghasilkan
apa-apa. Seberapa jauh maskapainya dapat mendorong sebelum penumpang yang tergusur
memakan labanya? \omterm{def:g12:randvar:binomial}{Distribusi binomial}
(\cref{thm:g12:randvar:binomial}) menjawabnya sampai ke angka desimalnya --- dan
mesin yang sama memeriksa lot pabrik, menetapkan harga undian, dan
menjaga penanggung asuransi tetap mampu bayar. Keputusan di bawah pengulangan: inilah
pekerjaan sehari-hari \omterm{def:g12:randvar:binomial}{distribusi binomial}.

\medskip
\noindent\textbf{Bagian I --- Kelancaran.}
\begin{enumerate}
  \item Berikan tabel \omterm{def:g12:randvar:rv}{distribusi} lengkap bagi cacah $X$ berisi
        angka enam pada tiga lemparan dadu.
  \item Hitunglah $\E(X)$ dan $V(X)$
        (\cref{def:g12:randvar:exp}).
  \item Sebuah kios memungut $1$ euro untuk tiga lemparan dan membayar $2$
        euro tiap angka enam yang diperoleh. Hitunglah harapan keuntungannya: adilkah
        permainannya?
  \item Latihan daftar periksa (\cref{met:g12:randvar:recognize}):
        apakah cacah kartu hati pada $5$ kartu yang diambil
        \emph{tanpa} pengembalian itu binomial? Dan dengan pengembalian?
        Berilah alasan.
  \item Untuk $X \sim \mathcal B(20,\ 0.3)$: berikan $\E(X)$ dan
        $\sigma(X)$, lalu $\E(S)$ dan $\sigma(S)$ bagi skor
        $S = 5X - 10$ (\cref{prop:g12:randvar:affine}).
\end{enumerate}

\medskip
\noindent\textbf{Bagian II --- Penerbangan yang kelebihan pesanan.}
Kursinya: $100$. Tiket yang terjual: $105$. Tiap pemegang tiket datang
dengan peluang $0.9$, secara \omterm{def:g12:condprob:indep}{saling bebas}; misalkan $X$ cacah orang yang
datang.
\begin{enumerate}[resume]
  \item Modelnya: berilah alasan bagi $X \sim \mathcal B(105,\ 0.9)$ dengan
        daftar periksanya --- lalu akuilah titik terlemah modelnya
        (benarkah ketidakhadirannya \omterm{def:g12:condprob:indep}{saling bebas}? pikirkanlah rombongan dan
        badai).
  \item Hitunglah $\E(X)$ dan $\sigma(X)$.
  \item Penggusuran terjadi ketika $X \geq 101$: tulislah
        $\P(X \geq 101)$ sebagai jumlah lima suku binomial yang
        gamblang.
  \item Hitunglah jumlahnya (dengan kalkulator): berapa bagian penerbangannya
        yang mengalami sekurang-kurangnya satu penumpang tergusur?
  \item Uangnya: kelima tiket tambahannya membawa
        $5 \times 200 = 1\,000$ euro; sedangkan tiap penumpang yang tergusur
        berbiaya $800$ euro sebagai ganti rugi. Hitunglah harapan
        cacah penumpang yang tergusur,
        $\sum_{k=101}^{105} (k - 100)\,\P(X = k)$, lalu harapan
        ganti ruginya, dan vonis atas kebijakannya.
  \item Perancangannya: pengaturnya menoleransi
        $\P(\text{penggusuran}) < 5\,\%$. Dengan menguji
        $n = 105, 106, 107, \dots$ tiket, carilah $n$ terbesar
        yang diperbolehkan.
  \item Jalan pintas loncengnya: untuk $n = 105$, hitunglah \omterm{pb:g11:stat:1}{skor-z}
        bagi ambang penggusurannya,
        $z = \frac{100.5 - \E(X)}{\sigma(X)}$, lalu bandingkan
        taksiran kasar ``kira-kira $2\sigma$, jadi kurang lebih $2$--$3\,\%$ pada
        ekor atasnya'' dengan jawaban eksakmu. (Kurva
        di balik jalan pintas ini adalah bintang bab
        berikutnya.)
\end{enumerate}

\medskip
\noindent\textbf{Bagian III --- Gerbang pabriknya.}
Sebuah lot suku cadang diterima jika contoh berisi $20$ buah memuat paling banyak
satu suku cadang yang cacat.
\begin{enumerate}[resume]
  \item Jika laju cacat yang sebenarnya adalah $2\,\%$ (lot yang jujur),
        hitunglah peluang penerimaannya.
  \item Jika lajunya $10\,\%$ (lot yang buruk), hitunglah lagi.
  \item Sebutkan kedua risiko prosedurnya (lot jujur yang
        ditolak; lot buruk yang diterima), bacalah nilainya dari
        pertanyaan 13--14, lalu katakan satu perubahan apa yang memperbaiki
        keduanya sekaligus --- dan dengan biaya apa.
  \item Di balik perbaikannya: tunjukkan bahwa \emph{frekuensi} cacat
        yang teramati, $\frac Xn$, pada contoh berukuran $n$ mempunyai
        \omterm{def:g11:stat:variance}{simpangan baku} $\sqrt{\frac{p(1-p)}{n}}$, lalu
        simpulkan bagaimana ketelitiannya berskala dengan ukuran contohnya (seorang
        kawan lama: $\frac{1}{\sqrt n}$ dari kelas 10).
\end{enumerate}

\medskip
\noindent\textbf{Bagian IV --- Permainan, undian, cadangan.}
\begin{enumerate}[resume]
  \item Permainan Saint Petersburg yang dibatasi pada
        \cref{exo:g12:randvar:8} berharapan bayaran $10$
        euro. Bandingkanlah dalam satu paragraf dengan paradoks tanpa
        batas yang dijumpai di \cref{pb:g11:prob:1}: apa tepatnya yang
        dijinakkan batas itu, dan harga berapa yang kini adil?
  \item Sebuah undian amal menjual $200$ tiket seharga $2$ euro;
        hadiahnya: satu keranjang $100$ euro dan dua keranjang $50$ euro.
        Hitunglah harapan keuntungan satu tiket dan marjin undiannya;
        lalu bandingkan dengan marjin bandar pada rolet Eropa
        sebesar $\frac{1}{37} \approx 2.7\,\%$. Mengapa undian itu
        lolos begitu saja?
  \item Seorang penanggung asuransi memegang $1\,000$ polis yang \omterm{def:g12:condprob:indep}{saling bebas}: peluang
        klaimnya $0.01$ masing-masing, besar klaimnya $10\,000$ euro, dan
        preminya $130$ euro. Hitunglah harapan laba tahunannya
        \emph{beserta} \omterm{def:g11:stat:variance}{simpangan baku} total
        klaimnya. Bandingkan kedua bilangan itu: apa yang dipaksakan
        pembandingan itu untuk dipegang penanggung asuransi yang sungguhan?
  \item Penutup --- \omterm{def:g12:randvar:binomial}{distribusi binomial} sebagai wasit keputusan: yaitu
        daftar periksa pengenalannya; pedoman $\E \pm \sigma$;
        peluang ekor sebagai harga risikonya; serta kedua peringatan
        yang tetap berlaku (\omterm{def:g12:condprob:indep}{kebebasan} itu klaim permodelan, dan
        ekor yang langka, bukan rata-ratanya, yang menyebabkan kebangkrutan). Masing-masing satu kalimat,
        dengan penunjuk: hukum bilangan besar dan kurva
        lonceng, pada bab berikutnya, melengkapi buku aturan sang wasit.

\end{enumerate}
\end{problem}
